\section{The Harmonic Cochain Problem}\label{sec:hcp-definition}

\begin{definition}[Public observation]\label{def:public-observation}
  A public HCP observation is a pair \((\alpha,b)\) consisting of a realized
  public classifier value \(\alpha\in\image(\pi)\) and an observed degree-one
  cochain \(b\in V\) in the fixed comparison coordinates.
\end{definition}

\begin{definition}[Hidden Hodge witness]\label{def:hidden-hodge-witness}
  A hidden Hodge witness for a public observation \((\alpha,b)\) is a hidden
  representative \(x\in A_\alpha\), satisfying \(\pi(x)=\alpha\). Its
  presentation data select the operative operator \(\Delta_x\). The witness
  and \(b\) determine the coordinates
  \[
    h_x(b):=P_{\mathcal H_x}b,
    \qquad
    e_x(b):=b-h_x(b).
  \]
  The Hodge decomposition in \cref{prop:hodge-decomposition} gives
  \[
    h_x(b)\in\ker(\Delta_x),\qquad
    e_x(b)\in
      \image(\delta_{0,x})\oplus\image(\delta_{1,x}^\ast),\qquad
    b=h_x(b)+e_x(b).
  \]
\end{definition}

\begin{definition}[Harmonic Cochain Problem]\label{def:hcp}
  The \emph{Harmonic Cochain Problem} (HCP) asks, given the public observation
  \((\alpha,b)\), to recover the operative harmonic component
  \[
    h=h_x(b)\in\mathcal H_x,
  \]
  equivalently the harmonic term in the orthogonal Hodge decomposition of
  \(b\) for the operative hidden witness \(x\in A_\alpha\). The witness is
  bound to the instance but is not part of the public input.
\end{definition}

\begin{definition}[HCP solution]\label{def:hcp-solution}
  Relative to an operative hidden witness \(x\), a solution to an
  HCP instance \((\alpha,b)\) is the cochain
  \(h'=h_x(b)=P_{\ker(\Delta_x)}b\). A public-data-only solver is
  \emph{uniformly exact}
  if the same output has this value for every hidden witness compatible with
  its public input.
\end{definition}

\subsection{Public Classes and Hidden Operators}\label{sec:mathematical-framework}

\begin{definition}[Hidden-operator public projection]\label{def:hidden-operator-public-projection}
  For a public class type \(P\), a public projection is a map
  \[
    \pi_\Delta:\{\text{hidden Hodge operators }\Delta\}\longrightarrow P.
  \]
  Two hidden operators are publicly equivalent when
  \(\pi_\Delta(\Delta)=\pi_\Delta(\Delta')\). The hidden operator fiber over
  \(\alpha\in P\) is
  \[
    \operatorname{Fib}_{\pi_\Delta}(\alpha)
    =\pi_\Delta^{-1}(\alpha)
    =\{\Delta:\pi_\Delta(\Delta)=\alpha\},
  \]
  and the hidden class of \(\Delta\) is
  \(\operatorname{Fib}_{\pi_\Delta}(\pi_\Delta(\Delta))\).
\end{definition}

\begin{definition}[Lossiness and kernel ambiguity]\label{def:kernel-ambiguity}
  The public projection \(\pi_\Delta\) is \emph{lossy} if there exist
  \(\Delta\ne\Delta'\) with
  \(\pi_\Delta(\Delta)=\pi_\Delta(\Delta')\). The two
  operators have the same harmonic kernel if
  \[
    \forall v,\quad \Delta v=0\Longleftrightarrow \Delta'v=0.
  \]
  The projection is \emph{kernel-ambiguous} if there are publicly equivalent
  operators \(\Delta,\Delta'\) whose harmonic kernels are not the same.
\end{definition}

\begin{definition}[Public-compatible Hodge ambiguity]\label{def:public-compatible-hodge-ambiguity}
  Public-compatible Hodge ambiguity means that there exist publicly equivalent
  operators \(\Delta,\Delta'\) and a cochain-facing observation \(b\) such that
  \[
    \Delta b=0
    \qquad\text{but}\qquad
    \Delta'b\ne 0.
  \]
\end{definition}

\begin{definition}[Typed HCP semantic realization]\label{def:typed-hcp-realization}
  Fix a multi-sorted algebraic signature \(\sigma\), a model category
  \(\mathcal C\) with terminal object and binary products, and, for each
  hidden representative \(x\in A\), a base
  interpretation \(M_x:\sigma\to\mathcal C_{\mathrm f}\) together with
  homotopical formula semantics over \(M_x\). Write
  \[
    \operatorname{Th}(M_x)
    =\{\varphi\text{ closed}:M_x\vDash\varphi\}
  \]
  for its closed first-order theory. Relative to the presentation data and
  coordinate charts fixed above, a typed HCP semantic realization consists
  additionally of maps
  \[
  \begin{aligned}
    x&\longmapsto\Delta_x\in\operatorname{End}(V),\\
    \kappa&:\{\text{closed theories}\}\longrightarrow P,\\
    \pi_\Delta&:\{\text{Hodge operators}\}\longrightarrow P,
  \end{aligned}
  \]
  satisfying the two factorization identities
  \[
    \pi(x)=\kappa(\operatorname{Th}(M_x)),
    \qquad
    \pi_\Delta(\Delta_x)=\pi(x).
  \]
  The Laplacian \(\Delta_x\) is presentation-level Hodge data extending the
  semantic core \(M_x\); it is not silently included among the symbols whose
  closed theory is asserted to be homotopy invariant.
\end{definition}

\begin{proposition}[Helfer-to-HCP classifier bridge]\label{prop:helfer-hcp-bridge}
  Let \(x,y\in A\) belong to a typed HCP semantic realization. If
  \(M_x\) and \(M_y\) are homotopy equivalent interpretations in the sense of
  Helfer, then
  \[
    \pi(x)=\pi(y)
    \qquad\text{and}\qquad
    \pi_\Delta(\Delta_x)=\pi_\Delta(\Delta_y).
  \]
  If, independently, there is a cochain \(b\) for which
  \[
    \Delta_xb=0,
    \qquad
    \Delta_yb\ne0,
  \]
  then \(\pi_\Delta\) has public-compatible Hodge ambiguity, and hence is
  kernel-ambiguous and lossy.
\end{proposition}

\begin{proof}
  The stated closed-test invariance premise gives, for every closed
  \(\sigma\)-formula \(\varphi\),
  \[
    M_x\vDash\varphi
    \quad\Longleftrightarrow\quad
    M_y\vDash\varphi.
  \]
  In the full first-order setting this premise is supplied by Helfer's
  homotopy-invariance theorem \cite[Theorem~12.3]{helfer2019fol}.
  Hence \(\operatorname{Th}(M_x)=\operatorname{Th}(M_y)\). Applying
  \(\kappa\) and then the two factorization identities gives
  \[
    \pi(x)=\pi(y),
    \qquad
    \pi_\Delta(\Delta_x)=\pi_\Delta(\Delta_y).
  \]
  The additional displayed harmonicity and non-harmonicity conditions now
  exhibit \((\Delta_x,\Delta_y,b)\) as a public-compatible Hodge-ambiguity
  witness. The ambiguity hierarchy below supplies kernel ambiguity and
  lossiness.
\end{proof}

Under closed-test invariance and classifier completeness,
\cref{prop:helfer-hcp-bridge} gives
\[
  \pi(x)=\pi(y)\quad\Longleftrightarrow\quad x\simeq y,
  \qquad
  A_{\pi(x)}=\{y\in A:y\simeq x\}.
\]
The implication from label equality uses completeness; the other uses
invariance. The fiber identity follows from the definition of \(A_{\pi(x)}\).
The family in \cref{prop:cellular-axis-ambiguity} has a complete classifier.
Completeness concerns the semantic class; it does not select its operative
embedded harmonic kernel. Kernel disagreement remains an independent condition. If
\(\Delta\)-harmonicity were included in the invariant first-order signature,
and the common coordinate \(b\) were
named compatibly in both interpretations, Helfer's theorem would instead rule
out the displayed disagreement.

Suppose an equivalence of hidden presentations induces degreewise linear
isometric equivalences
\(U_0,U_1,U_2\) satisfying
\[
  U_1\delta_{0,x}=\delta_{0,y}U_0,
  \qquad
  U_2\delta_{1,x}=\delta_{1,y}U_1.
\]
In the fixed coordinate charts, \(U_1\) is an isometry of \(V\). The
isometries transport adjoints, giving
\[
  \Delta_yU_1=U_1\Delta_x,
  \qquad
  U_1(\ker\Delta_x)=\ker\Delta_y.
\]
The corresponding orthogonal projections commute with transport:
\[
  U_1P_{\ker\Delta_x}b
  =
  P_{\ker\Delta_y}U_1b.
\]
Thus a transported observation preserves both the operative target and
harmonicity:
\[
  \Delta_xb=0
  \quad\Longleftrightarrow\quad
  \Delta_y(U_1b)=0.
\]
The ambiguity witness in \cref{prop:helfer-hcp-bridge} uses a fixed public
coordinate \(b\in V\). Under the displayed conjugacy, its separation condition is
equivalently
\[
  b\in\ker\Delta_x,
  \qquad
  b\notin U_1(\ker\Delta_x).
\]
It requires a nonstabilization witness in addition to the equivalence, with
the private transport \(U_1\) neither revealed nor applied.

\begin{proposition}[Ambiguity hierarchy]\label{prop:ambiguity-hierarchy}
  Public-compatible Hodge ambiguity implies kernel ambiguity, and kernel
  ambiguity implies that the public projection is lossy.
\end{proposition}

\begin{proof}
  If \(b\) is harmonic for \(\Delta\) but not for \(\Delta'\), then the two
  operators cannot have the same harmonic kernel. This gives kernel ambiguity.
  If publicly equivalent operators have different kernels, they cannot be the
  same operator, since equality of operators implies equality of kernels.
  Therefore the public projection identifies two distinct hidden operators.
\end{proof}

\subsection{Computational Scope}\label{sec:computational-scope}

These exact-recovery comparisons do not by themselves give an asymptotic
CVP-hardness reduction. Such a comparison requires encoded instance families,
an input-size measure, an algorithm class and a success criterion. A
distributional security claim also requires a sampling model and a declared
adversary view.
